% ============================================================================
% Fichier  : partie5/P05_C15_cycle_reponse.tex
% Titre    : Le Cycle de Réponse à Incident
% Auteur   : MINKA MI NGUIDJOÏ Thierry Emmanuel
% ============================================================================
% Sources : ISO/IEC 27035 ; NIST SP 800-61r2. Opération MAYO : fictif.
% ============================================================================

\chapter{Le Cycle de Réponse à Incident}
\label{chap:cycle}
\chaptermark{Le cycle de réponse}

\epigraph{%
    « La réponse se prépare à froid~;\\
    en pleine crise, on applique, on n'improvise pas. »%
}{Principe du cours}

\niveauChapitre{Fondamental}{Chapitre~\ref{chap:normes}.}

\begin{boxobjectifs}
\begin{itemize}
    \item Décrire le cycle de réponse (ISO 27035, NIST 800-61r2).
    \item Classer un incident par gravité et fixer les délais.
    \item Comprendre le rôle de la préparation.
    \item Distinguer confinement, éradication, rétablissement.
    \item Articuler réponse et forensique sans les opposer.
\end{itemize}
\end{boxobjectifs}

\bigskip

% ============================================================================
\section*{Question directrice : le 6 mars à 7 h 15, que faut-il faire, et dans quel ordre~?}
% ============================================================================

Devant l'incident, l'urgence pousse à tout faire en même temps. Le cycle de
réponse impose un ordre~: qualifier, décider, contenir, éradiquer, rétablir,
apprendre. Ce chapitre le pose, avec les délais par gravité.

\medskip
\textbf{Plan du chapitre.}
\begin{enumerate}
    \item Le cycle (section~\ref{sec:c15-cycle}).
    \item Classer par gravité (section~\ref{sec:c15-gravite}).
    \item La préparation (section~\ref{sec:c15-preparation}).
    \item Réponse et forensique (section~\ref{sec:c15-forensique}).
\end{enumerate}

\separateur

% ============================================================================
\section{Le Cycle}
\label{sec:c15-cycle}
% ============================================================================

\begin{tableaucours}{L{3.6cm}L{8.6cm}}{%
    \tableauHeader{Étape} &
    \tableauHeader{Ce qu'on y fait}
}
Préparer & Politique, rôles, playbooks, outils, astreinte \\
Détecter et qualifier & Confirmer l'incident, évaluer la gravité \\
Contenir & Limiter la propagation sans détruire la preuve \\
Éradiquer & Supprimer la cause et la présence de l'attaquant \\
Rétablir & Restaurer le service à partir de sauvegardes vérifiées \\
Apprendre & Retour d'expérience, amélioration (Partie~VI) \\
\end{tableaucours}
\vspace{-0.8em}
\begin{center}{\small\itshape Tableau~15.1. Le cycle de réponse~\cite{iso27035,nist80061r2}}\end{center}

% ============================================================================
\section{Classer par Gravité}
\label{sec:c15-gravite}
% ============================================================================

\begin{tableaucours}{L{2.4cm}L{5.2cm}L{4.2cm}}{%
    \tableauHeader{Gravité} &
    \tableauHeader{Exemple} &
    \tableauHeader{Délai de prise en charge}
}
Critique & Rançongiciel, arrêt de production & Moins d'une heure, astreinte \\
Élevée & Compromission d'un compte à privilèges & Moins de 4 heures \\
Moyenne & Poste infecté isolé & Moins de 24 heures \\
Faible & Tentative bloquée & Moins de 72 heures \\
\end{tableaucours}
\vspace{-0.8em}
\begin{center}{\small\itshape Tableau~15.2. Grille de gravité et délais (LiZbiZ)}\end{center}

% ============================================================================
\section{La Préparation}
\label{sec:c15-preparation}
% ============================================================================

La préparation est la phase la plus déterminante~: la qualité de la réponse
se joue avant l'incident. Elle comprend les rôles, les coordonnées, les
playbooks (rançongiciel, hameçonnage, fuite de données), les sauvegardes
vérifiées et les droits d'accès d'urgence. Un incident révèle toujours l'état
de la préparation.

\begin{boxalert}
\textbf{Erreur fréquente.} Rétablir avant d'éradiquer.

\smallskip
\textbf{Règle~:} restaurer un service alors que l'attaquant est toujours
présent le ramène aussitôt. On éradique (et on vérifie l'éradication) avant
de rétablir.
\end{boxalert}

% ============================================================================
\section{Réponse et Forensique}
\label{sec:c15-forensique}
% ============================================================================

\begin{boxscenario}{Opération MAYO. L'ordre tenu}
\textbf{Constatation.} Le 6~mars, après qualification (critique), Mme~MOUAFO
décide~: isoler les systèmes atteints (confinement), capturer les preuves
avant toute remise en service, éradiquer le compte \texttt{svc\_maj2} et le
service \texttt{WinUpd32}, puis rétablir depuis des sauvegardes vérifiées hors
ligne.

\medskip
\textbf{Analyse} (confiance élevée). Le confinement a précédé le
rétablissement~; la capture des preuves (Partie~IV) a été insérée avant
l'éradication. Réponse et forensique ne s'opposent pas~: on contient, on fige,
puis on nettoie.
\end{boxscenario}

\begin{boxCRO}
\textbf{Contenir sans détruire la preuve.} Isoler une machine du réseau
préserve la preuve (on ne l'éteint pas) et limite la propagation~:
\cro{+}{+}{+}. Éteindre brutalement contiendrait aussi, mais détruirait la
mémoire~: \cro{+}{-}{-}. On préfère isoler.
\end{boxCRO}

\begin{boxmetier}{Responsable de la réponse à incident}
\textbf{Tenir l'ordre sous pression.} Le RRI résiste à la pression de
« tout rallumer tout de suite »~: il fait tenir l'ordre confiner, figer,
éradiquer, rétablir. C'est la discipline qui évite la rechute.
\end{boxmetier}

\separateur

\begin{boxresume}
\textbf{Ce qu'il faut retenir de ce chapitre~:}
\begin{itemize}
    \item \textbf{Préparer, qualifier, contenir, éradiquer, rétablir,
          apprendre.}
    \item \textbf{La gravité} fixe les délais.
    \item \textbf{La préparation} décide de la qualité de la réponse.
    \item \textbf{Éradiquer avant de rétablir}, sous peine de rechute.
    \item \textbf{Contenir en isolant}, pour préserver la preuve.
\end{itemize}
\end{boxresume}

\begin{boxinfo}
\textbf{Points de contrôle de l'élève.} Je sais~: décrire le cycle~; classer
par gravité~; expliquer le rôle de la préparation~; ordonner éradication et
rétablissement~; articuler réponse et forensique.
\end{boxinfo}

\bigskip

\begin{boxexo}[title={\ding{45}\quad Exercice~15.1. Le cycle \niveaubase}]
Citez les six étapes du cycle de réponse.
\end{boxexo}

\begin{boxexo}[title={\ding{45}\quad Exercice~15.2. Gravité \niveaubase}]
À quelle gravité classer un rançongiciel qui arrête la production~?
\end{boxexo}

\begin{boxexo}[title={\ding{45}\quad Exercice~15.3. Ordre \niveaubase}]
Pourquoi ne pas rétablir avant d'éradiquer~?
\end{boxexo}

\begin{boxexo}[title={\ding{45}\quad Exercice~15.4. Préparation \niveaumoyen}]
Pourquoi la préparation est-elle la phase la plus déterminante~?
\end{boxexo}

\begin{boxexo}[title={\ding{45}\quad Exercice~15.5. Isoler \niveaumoyen}]
Pourquoi isoler plutôt qu'éteindre~? (Analyse CRO)
\end{boxexo}

\begin{boxexo}[title={\ding{45}\quad Exercice~15.6. MAYO \niveaumoyen}]
Reconstituez l'ordre des décisions de Mme~MOUAFO le 6~mars.
\end{boxexo}

\begin{boxexo}[title={\ding{45}\quad Exercice~15.7. Figer \niveaumoyen}]
Où insère-t-on la capture des preuves dans le cycle~?
\end{boxexo}

\begin{boxexo}[title={\ding{45}\quad Exercice~15.8. Sauvegardes \niveaumoyen}]
Pourquoi rétablir à partir de sauvegardes vérifiées hors ligne~?
\end{boxexo}

\begin{boxexo}[title={\ding{45}\quad Exercice~15.9. Playbook \niveauexpert}]
Rédigez le playbook rançongiciel de LiZbiZ, étape par étape, avec les
décisions.
\end{boxexo}

\begin{boxexo}[title={\ding{45}\quad Exercice~15.10. Grille \niveauexpert}]
Construisez la grille de gravité de LiZbiZ avec délais, astreinte et
escalades.
\end{boxexo}

\bigskip

\begin{boxinfo}
\textbf{Pour aller plus loin.}
\begin{itemize}
    \item ISO/IEC 27035~\cite{iso27035}~; NIST SP~800-61r2~\cite{nist80061r2}.
\end{itemize}
\end{boxinfo}

\begin{boxinfo}
\textbf{Transition.} Répondre suppose une équipe organisée~: le
Chapitre~\ref{chap:csirt} traite du CSIRT et du CERT.
\end{boxinfo}
